% Source: https://artofproblemsolving.com/community/c3h1521268_2014_mathcounts_national_sprint_problems
% Stable identity: aops-topic-1521268-q13
Let the corner numbers have sum $c$ and let each side sum be $T$. Adding the three side sums counts each corner twice and each middle number once:
\[
3T=(1+2+\cdots+6)+c=21+c .
\]
So $c$ is a multiple of $3$ between $1+2+3=6$ and $4+5+6=15$: $c\in\{6,9,12,15\}$, with $T=9,10,11,12$. Once the corners are placed, each middle number is forced: it is $T$ minus its two neighbouring corners. Also any arrangement of three given corner numbers can be turned into any other by a rotation or a reflection. So for each $c$ we only need to find which sets of corners work.
\begin{itemize}
\item $c=6$, $T=9$: corners $\{1,2,3\}$, middles $9-1-2=6$, $9-2-3=4$, $9-3-1=5$. Valid.
\item $c=9$, $T=10$: corner sets $\{1,2,6\}$ (middle $10-1-2=7$, too big), $\{1,3,5\}$ (middles $6,2,4$, valid), $\{2,3,4\}$ (middle $10-2-4=4$ repeats a corner). One valid.
\item $c=12$, $T=11$: corner sets $\{1,5,6\}$ (middle $11-1-5=5$ repeats), $\{2,4,6\}$ (middles $5,1,3$, valid), $\{3,4,5\}$ (middle $11-3-4=4$ repeats). One valid.
\item $c=15$, $T=12$: corners $\{4,5,6\}$, middles $12-4-5=3$, $12-5-6=1$, $12-6-4=2$. Valid.
\end{itemize}
Hence there are $\boxed{4}$ different solutions.
